% !TeX root = ../main.tex
\chapter{Analiza sentymentu: podstawy i przegląd literatury}
\label{chap:sentiment}

% FIRST DRAFT. Claims about the literature follow the abstracts of the cited
% works; check each against the full text before submission.

Zanim opiszemy system, trzeba ustalić, co właściwie mierzy.
Słowa \enquote{opinia}, \enquote{sentyment}, \enquote{wydźwięk} i \enquote{nastrój} bywają w literaturze i w mowie potocznej używane zamiennie, a w tej pracy każde z nich oznacza coś innego.
Rozdział zaczyna się więc od definicji, następnie przegląda metody automatycznej klasyfikacji wydźwięku, od słownikowych po modele typu transformer, i kończy się omówieniem badań, które śledzą opinie w czasie, wokół konkretnych wydarzeń.

\section{Pojęcia i przyjęta terminologia}
\label{sec:sentiment-terms}

Pang i Lee w przeglądzie z 2008~r. traktują eksplorację opinii (ang. \textit{opinion mining}) i analizę sentymentu (ang. \textit{sentiment analysis}) jako dwie nazwy jednej dziedziny: automatycznego rozpoznawania, co autor tekstu sądzi o jego przedmiocie~\cite{pang2008opinion}.
Liu proponuje bardziej formalną definicję, w której opinia jest piątką: obiekt, którego dotyczy, jego aspekt, wydźwięk wypowiedzi, autor i czas jej wyrażenia~\cite{liu2015sentiment}.
Ta definicja dobrze pasuje do naszego problemu, bo dwa jej elementy, autor i czas, są w systemie pierwszoplanowe.
Autora reprezentujemy wyłącznie pseudonimem, a czas wyznacza miejsce opinii w odtwarzanym strumieniu.

Liu wyróżnia też trzy poziomy analizy: całego dokumentu, pojedynczego zdania i aspektu~\cite{liu2015sentiment}.
Analiza aspektowa odpowiedziałaby na pytanie, co dokładnie się podoba, na przykład fabuła, a co nie, na przykład wydajność.
My pracujemy na poziomie dokumentu, czyli całego komentarza albo całej recenzji.
Recenzja Steam ma średnio 248 znaków, a komentarz YouTube 123 znaki, więc w większości przypadków dokument ma długość jednego lub dwóch zdań i poziom dokumentu niewiele różni się od poziomu zdania.

W pracy używamy następujących pojęć.
\emph{Wydźwięk} (sentyment) to klasa przypisana jednej opinii: pozytywna, neutralna albo negatywna.
Klasa neutralna obejmuje zarówno wypowiedzi bez oceny, na przykład pytania o datę premiery, jak i wypowiedzi, w których pochwała i krytyka się równoważą.
\emph{Nastrój} to stan zbiorowości, a nie pojedynczej osoby.
Mierzymy go przez agregację wydźwięku wielu opinii z określonego okresu, źródła lub okna wydarzenia.
Do agregacji służy \emph{wskaźnik wydźwięku netto}, zdefiniowany jako
\begin{equation}
  N = \frac{n_{\mathrm{pos}} - n_{\mathrm{neg}}}{n_{\mathrm{pos}} + n_{\mathrm{neu}} + n_{\mathrm{neg}}},
  \label{eq:net-sentiment}
\end{equation}
gdzie $n_{\mathrm{pos}}$, $n_{\mathrm{neu}}$ i $n_{\mathrm{neg}}$ to liczby opinii w poszczególnych klasach.
Wskaźnik przyjmuje wartości od $-1$ do $1$ i ma tę zaletę, że da się go policzyć dla dowolnej sumy grup, dodając liczniki i dzieląc raz na końcu.
Średnia wskaźników policzonych osobno dla grup o różnej liczebności dałaby inny, mylący wynik, dlatego warstwa raportowa przechowuje liczniki, a nie gotowe udziały.

Od wydźwięku odróżniamy \emph{emocję}.
Klasyfikacja emocji przypisuje tekstowi jedną z wielu kategorii, takich jak radość, złość czy zaskoczenie; zbiór GoEmotions rozróżnia ich 27, plus klasę neutralną~\cite{demszky2020goemotions}.
Emocje nie są przedmiotem tej pracy, choć zbiór GoEmotions wykorzystaliśmy pośrednio, jako część danych treningowych po sprowadzeniu jego etykiet do trzech klas wydźwięku.

Dwa ostatnie pojęcia dotyczą oceny modelu.
\emph{Zbiór wzorcowy} to próbka opinii z naszych danych, której wydźwięk oznaczyli ludzie według spisanych wytycznych.
\emph{Słaba etykieta} to sygnał o wydźwięku, który nie pochodzi od oceniającego, ale da się go uzyskać bez dodatkowej pracy.
W recenzjach Steam taką etykietą jest deklaracja autora, czy poleca grę.
Etykieta ta jest słaba, bo recenzja polecająca grę może zawierać ostre uwagi, ale ma dwie zalety: jest dostępna dla każdej z ponad \num{250000} recenzji i nie zależy od nas.

\section{Metody klasyfikacji wydźwięku}
\label{sec:sentiment-methods}

Metody klasyfikacji wydźwięku rozwijały się w trzech falach.
Każda z nich jest wciąż używana, bo każda ma inny stosunek kosztu do jakości.

\subsection{Metody słownikowe}
\label{sec:sentiment-lexicon}

Najstarsze podejście korzysta ze słownika, w którym każdemu słowu przypisano wydźwięk i jego siłę.
Wydźwięk tekstu to suma wydźwięków słów, poprawiona regułami dla zaprzeczeń (\enquote{nie najlepszy}) i wzmocnień (\enquote{bardzo dobry}).
Dobrze widać to na modelu VADER, który Hutto i Gilbert przygotowali z myślą o mediach społecznościowych.
Bierze on pod uwagę wielkie litery, wykrzykniki, emotikony i slang, a na krótkich wpisach z Twittera osiągnął w ocenie autorów lepszą trafność niż pojedynczy ludzcy oceniający~\cite{hutto2014vader}.

Metody słownikowe są tanie, szybkie i łatwe do wyjaśnienia, bo zawsze wiadomo, które słowa przesądziły o wyniku.
Ich słabością jest zależność od języka i dziedziny.
Słowo \enquote{brutalny} w recenzji gry akcji bywa pochwałą, a słowniki dla języków innych niż angielski są mniejsze albo nie ma ich wcale.
Przy danych w kilkudziesięciu językach musielibyśmy utrzymywać kilkadziesiąt słowników.
W projekcie rolę podejścia słownikowego pełniła prosta reguła słów kluczowych, którą wykorzystywaliśmy jako punkt odniesienia; na zbiorze wzorcowym osiągnęła dokładność \num{0.417}.

\subsection{Klasyczne uczenie maszynowe}
\label{sec:sentiment-classic-ml}

Pang, Lee i Vaithyanathan potraktowali wydźwięk recenzji filmowych jako zadanie klasyfikacji tekstu i porównali naiwny klasyfikator bayesowski, model maksymalnej entropii i maszynę wektorów nośnych~\cite{pang2002thumbs}.
Modele uczone na cechach w postaci obecności słów okazały się lepsze od reguł ułożonych przez ludzi, ale autorzy zauważyli też, że wydźwięk jest trudniejszy do klasyfikacji niż temat tekstu.
Zdanie \enquote{fabuła miała wszystko, żeby się udać} zawiera same pozytywne słowa, a wyraża rozczarowanie.

Klasyczne modele potrzebują oznaczonych danych, a ręczne oznaczanie jest drogie.
Jednym z obejść jest nadzór odległy (ang. \textit{distant supervision}): Go, Bhayani i Huang oznaczyli wpisy z Twittera automatycznie, na podstawie emotikonów, i na tak zebranym zbiorze, znanym później jako Sentiment140, wytrenowali klasyfikatory wydźwięku~\cite{go2009twitter}.
Ten sam pomysł, czyli wykorzystanie sygnału, który autor tekstu zostawił przy okazji, wraca w naszej pracy w postaci rekomendacji autorów recenzji Steam.

\subsection{Sieci neuronowe i modele typu transformer}
\label{sec:sentiment-transformers}

Kolejnym krokiem były sieci neuronowe, które uczą się reprezentacji tekstu zamiast korzystać z ręcznie zaprojektowanych cech.
Socher i współautorzy zbudowali zbiór Stanford Sentiment Treebank, w którym wydźwięk oznaczono nie tylko dla całych zdań, ale i dla ich fragmentów, oraz rekurencyjną sieć, która składa wydźwięk zdania z wydźwięków jego części~\cite{socher2013recursive}.
Binarna wersja tego zbioru, SST-2, stała się jednym ze standardowych testów modeli językowych.

Przełom przyniosła architektura transformer, oparta wyłącznie na mechanizmie uwagi, bez warstw rekurencyjnych~\cite{vaswani2017attention}.
Model BERT pokazał, że transformer wstępnie wytrenowany na dużym korpusie bez etykiet, z zadaniem odgadywania zamaskowanych słów, można potem dostosować do konkretnego zadania przez dostrajanie (ang. \textit{fine-tuning}), dodając tylko jedną warstwę wyjściową~\cite{devlin2019bert}.
Ten schemat, czyli drogie trenowanie ogólne raz, a potem tanie dostrajanie do własnego problemu, jest dokładnie tym, czego potrzebowaliśmy.
Nie mieliśmy ani danych, ani budżetu, żeby trenować model od zera.

Dla systemu, który ma działać tanio, liczy się również rozmiar modelu.
DistilBERT, uzyskany przez destylację wiedzy z modelu BERT, jest o \qty{40}{\percent} mniejszy i o \qty{60}{\percent} szybszy, a zachowuje około \qty{97}{\percent} jego zdolności rozumienia języka~\cite{sanh2019distilbert}.
To dlatego był naszym pierwszym wyborem.
Dla danych wielojęzycznych odpowiednikiem jest XLM-RoBERTa (XLM-R), model z rodziny BERT wstępnie wytrenowany na ponad dwóch terabajtach przefiltrowanego tekstu w stu językach~\cite{conneau2020unsupervised}.
Ważną własnością takiego modelu jest transfer międzyjęzykowy: model dostrojony na danych w kilku językach klasyfikuje też teksty w językach, których nie było w danych dostrajających, choć zwykle z niższą trafnością.
Z tej własności korzystamy na przykład dla recenzji po polsku, dla których nie mieliśmy żadnych oznaczonych danych.

Gotowe implementacje tych modeli, wraz z wagami wstępnie wytrenowanych wersji, udostępnia biblioteka Transformers firmy Hugging Face~\cite{wolf2020transformers}, z której korzystaliśmy przy trenowaniu i serwowaniu modeli.
Z repozytorium modeli Hugging Face pochodzi też nasz model bazowy, \code{xlm-roberta-base}.

\subsection{Zbiory danych i specyfika mediów społecznościowych}
\label{sec:sentiment-social-media}

Teksty z mediów społecznościowych różnią się od recenzji filmowych, na których powstawały pierwsze metody.
Są krótkie, pełne skrótów, emotikonów, wzmianek o innych użytkownikach i adresów stron, często ironiczne.
Dlatego powstały dla nich osobne zbiory danych.
Zbiór TweetEval ujednolicił siedem zadań klasyfikacji wpisów, w tym wydźwięk, ironię i emocje, pod wspólnym protokołem oceny~\cite{barbieri2020tweeteval}.
XLM-T rozszerzył to podejście na wiele języków i udostępnił zbiór wpisów z oznaczonym wydźwiękiem w ośmiu językach~\cite{barbieri2022xlmt}.
Wielojęzyczny zbiór opinii z trzema klasami wydźwięku udostępnia też zespół clapAI w serwisie Hugging Face~\cite{clapaimultilingualsentiment}.

Trzy z tych zbiorów (TweetEval, XLM-T i clapAI), uzupełnione o GoEmotions, złożyły się na dane, na których dostrajaliśmy model XLM-R.
Żaden z nich nie dotyczy gier, więc różnica dziedzin pozostaje.
Zmniejszaliśmy ją, dodając do danych treningowych oznaczone komentarze z naszego zbioru, a jakość mierzyliśmy wyłącznie na komentarzach z naszego zbioru, odłożonych przed trenowaniem.

Na koniec warto wspomnieć o podejściu, które zyskało popularność w ostatnich latach: klasyfikacji wydźwięku przez duże modele językowe (\ac{llm}) sterowane poleceniem, bez dostrajania.
\todo{Źródło o klasyfikacji wydźwięku przez LLM bez dostrajania: find-sources.}
Nie wybraliśmy go z dwóch powodów.
Koszt wywołania dużego modelu dla każdej z ćwierci miliona opinii przekraczał nasz budżet, a odpowiedzi modelu zdalnego, którego wersji nie kontrolujemy, trudno uczynić powtarzalnymi.

\section{Analiza opinii w czasie wokół wydarzeń}
\label{sec:sentiment-time}

Pojedyncza ocena wydźwięku rzadko jest celem sama w sobie.
Ciekawsze jest to, jak nastroje zmieniają się w czasie i czy te zmiany mówią coś o świecie poza siecią.

O'Connor i współautorzy porównali wydźwięk wpisów z Twittera z sondażami zaufania konsumentów i poparcia dla prezydenta Stanów Zjednoczonych.
Po wygładzeniu szeregów czasowych korelacja w niektórych zestawieniach sięgała \num{0.8}, choć autorzy zastrzegli, że prosta metoda słownikowa jest daleka od precyzyjnej~\cite{oconnor2010tweets}.
Asur i Huberman pokazali, że sama liczba wpisów o filmie przed premierą pozwala przewidzieć jego wpływy kasowe lepiej niż rynek prognoz, a wydźwięk wpisów po premierze dodatkowo poprawia prognozę~\cite{asur2010predicting}.
Premiera filmu to wydarzenie tego samego rodzaju co premiery, które śledzimy, tyle że nas interesuje sam nastrój, a nie to, co mówi o sprzedaży.

Najbliżej naszego problemu leży praca Thelwalla, Buckleya i Paltoglou, którzy przyjrzeli się wydźwiękowi wpisów wokół głośnych wydarzeń.
Ich wynik był przekorny: popularne wydarzenia zwykle wiązały się ze wzrostem siły wydźwięku negatywnego, nawet gdy samo wydarzenie nie było złą wiadomością~\cite{thelwall2011sentiment}.
Wynika z tego praktyczny wniosek dla naszej analizy.
Porównywanie okna wydarzenia z resztą okresu ma sens tylko wtedy, gdy punkt odniesienia pochodzi z tego samego źródła.
Komentarze pod zwiastunem i recenzje kupujących różnią się już na wejściu, niezależnie od wydarzenia, dlatego efekt wydarzenia liczymy względem tła z tego samego serwisu (podrozdział~\ref{sec:impl-reporting}).

Recenzje gier w sklepie Steam badali Lin i współautorzy, traktując je jako źródło informacji zwrotnej dla twórców~\cite{lin2019empirical}.
Dla nas Steam ma jeszcze jedną zaletę, opisaną w poprzednim podrozdziale: autor recenzji sam deklaruje, czy poleca grę.
Pozwala to sprawdzić model na całym zbiorze, a nie tylko na kilkuset ręcznie oznaczonych opiniach.

Przegląd prac prowadzi do trzech obserwacji, które ukształtowały projekt systemu.
Po pierwsze, wiarygodny przebieg nastrojów wymaga danych z więcej niż jednego źródła, bo każde źródło ma inny profil użytkowników i inny zwyczajowy ton.
Po drugie, wynik zależy od metody klasyfikacji na tyle mocno, że trzeba umieć go policzyć ponownie, gdy metoda się zmieni.
To oznacza przechowywanie surowych danych i możliwość odtworzenia całego przetwarzania.
Po trzecie, przy danych wielojęzycznych model wielojęzyczny jest koniecznością, a nie ulepszeniem, co boleśnie sprawdziliśmy na pierwszym, angielskim modelu.
Jak te obserwacje przełożyły się na wybór konkretnych technologii, opisuje następny rozdział.
